\section{Introducción}
Las personas con diabetes necesitan seguir la evolución de su glucosa y relacionarla con lo que ocurre durante el día. Los monitores continuos de glucosa (CGM) permiten observar variaciones y tendencias, pero interpretar su relación con las comidas, el estrés o la actividad física suele requerir registros adicionales y evaluación profesional.

Este proyecto propone complementar ese seguimiento con un sistema formado por un parche que capta biomarcadores del sudor, un lector electrónico reutilizable, una pulsera y una aplicación móvil. El lector obtiene mediciones del parche y la pulsera registra frecuencia cardíaca, movimiento y temperatura de piel. La aplicación reúne ambas fuentes para mostrar el historial, explorar factores asociados a las variaciones y evaluar predicciones de corto plazo y de horizonte extendido.

La propuesta distingue las mediciones en sudor de las estimaciones de glucosa sanguínea. También distingue una asociación con hábitos o señales fisiológicas de una causa clínica. La calidad de cada salida deberá contrastarse con una referencia adecuada antes de habilitarla para el uso previsto. El sistema no emitirá recomendaciones automáticas de tratamiento.

El plan prevé un prototipo integrado en el mes 7 y un producto mínimo viable (MVP) técnico en el mes 12. La preparación clínica y regulatoria comienza durante el desarrollo y su cierre se proyecta hacia el mes 18. Las ventas se estiman desde el mes 19, sujetas a las autorizaciones correspondientes. Los beneficios esperados y las funciones predictivas se presentan como objetivos a evaluar, no como resultados ya demostrados.
